% !TeX root = ../sections/02_exercises.tex
\subsection{1-A-5.}
\begin{exercise}
	数列 $\{a_n\}$ と $\{b_n\}$ はそれぞれ
	$\alpha\in\mathbb{R}$ および $\beta\in\mathbb{R}$ に収束するとき，
	$\lambda,\mu\in\mathbb{R}$ に対して数列
	\[
		\{\lambda a_n+\mu b_n\}
	\]
	は収束して
	\[
		\lim_{n\to\infty}(\lambda a_n+\mu b_n)
		=
		\lambda\alpha+\mu\beta
	\]
	であることを証明せよ。
\end{exercise}
\begin{background}
	この問題では，収束する数列の線形結合も収束し，
	その極限は極限値の線形結合になることを示す。

	\begin{definition}[数列の収束]
		数列 $\{a_n\}$ が実数 $\alpha$ に収束するとは，
		任意の $\epsilon>0$ に対して，ある自然数 $N$ が存在して，
		すべての $n\geq N$ に対して
		\[
			|a_n-\alpha|<\epsilon
		\]
		が成り立つことをいう。
	\end{definition}

	\begin{remark}
		今回示したいことは
		\[
			\lambda a_n+\mu b_n
			\to
			\lambda\alpha+\mu\beta
		\]
		である。

		したがって，極限の定義より，
		任意の $\epsilon>0$ に対して，十分大きな $n$ で
		\[
			|(\lambda a_n+\mu b_n)-(\lambda\alpha+\mu\beta)|<\epsilon
		\]
		を示せばよい。
	\end{remark}

	\begin{lemma}[三角不等式]
		任意の実数 $x,y$ に対して
		\[
			|x+y|\leq |x|+|y|
		\]
		が成り立つ。
	\end{lemma}

	\begin{lemma}[定数倍と絶対値]
		任意の実数 $c,x$ に対して
		\[
			|cx|=|c||x|
		\]
		が成り立つ。
	\end{lemma}

	\begin{remark}
		差を整理すると，
		\[
			(\lambda a_n+\mu b_n)-(\lambda\alpha+\mu\beta)
			=
			\lambda(a_n-\alpha)+\mu(b_n-\beta)
		\]
		である。

		したがって，三角不等式より
		\[
			|(\lambda a_n+\mu b_n)-(\lambda\alpha+\mu\beta)|
			\leq
			|\lambda||a_n-\alpha|
			+
			|\mu||b_n-\beta|
		\]
		と評価できる。
	\end{remark}
\end{background}
\begin{strategy}
	示したいことは
	\[
		\lim_{n\to\infty}(\lambda a_n+\mu b_n)
		=
		\lambda\alpha+\mu\beta
	\]
	である。

	まず，極限の定義に戻る。
	任意の $\epsilon>0$ に対して，十分大きな $n$ で
	\[
		|(\lambda a_n+\mu b_n)-(\lambda\alpha+\mu\beta)|<\epsilon
	\]
	を示せばよい。

	差を変形すると，
	\[
		(\lambda a_n+\mu b_n)-(\lambda\alpha+\mu\beta)
		=
		\lambda(a_n-\alpha)+\mu(b_n-\beta)
	\]
	である。

	したがって，三角不等式を用いて
	\[
		|(\lambda a_n+\mu b_n)-(\lambda\alpha+\mu\beta)|
		\leq
		|\lambda||a_n-\alpha|
		+
		|\mu||b_n-\beta|
	\]
	と評価する。

	ここで，$a_n\to\alpha$ であるから，
	$n$ を十分大きくすれば $|a_n-\alpha|$ を小さくできる。
	また，$b_n\to\beta$ であるから，
	$n$ を十分大きくすれば $|b_n-\beta|$ も小さくできる。

	具体的には，
	\[
		|a_n-\alpha|<\frac{\epsilon}{2(|\lambda|+1)}
	\]
	かつ
	\[
		|b_n-\beta|<\frac{\epsilon}{2(|\mu|+1)}
	\]
	となるように $n$ を十分大きく取る。

	このように選ぶと，
	\[
		|\lambda||a_n-\alpha|
		<
		\frac{|\lambda|}{|\lambda|+1}\cdot\frac{\epsilon}{2}
		\leq
		\frac{\epsilon}{2}
	\]
	であり，同様に
	\[
		|\mu||b_n-\beta|
		<
		\frac{\epsilon}{2}
	\]
	となる。

	したがって，十分大きな $n$ に対して
	\[
		|(\lambda a_n+\mu b_n)-(\lambda\alpha+\mu\beta)|
		<
		\epsilon
	\]
	が成り立つ。

	つまり，線形結合の誤差を
	\[
		a_n-\alpha
		\quad\text{と}\quad
		b_n-\beta
	\]
	の誤差に分解し，それぞれを十分小さくすることが方針である。
\end{strategy}